% !TEX root = ../main.tex
\chapter{Newtonian gravitational interaction (Newtonsche Gravitationswechselwirkung)}\label{ch:c11}
\begin{paracol}{2}

\begin{sources}
\begin{tabularx}{\linewidth}{@{}l l L@{}}
\toprule
\textbf{Video} & \textbf{Length} & \textbf{Content}\\
\midrule
\seg{11}{1}{Lecture 11: Newtonian gravitational interaction} & 1:07:47 & Newton's achievement; the interaction fields are produced by all particles (active versus passive charge), the coupled problem and the instantaneous Newtonian field equations;
the covector field $W_\alpha$; the three Newtonian gravitational field equations; excursion on simple connectedness; consequences in a flat chart: the Poisson equation; the general solution and the potential of a spherically symmetric body; equivalence of inertial and gravitational mass and the start of the two-particle solution.\\
\bottomrule
\end{tabularx}

\smallskip
\textbf{Transcript} \file{materials/transcripts/C11-V01.txt} (German YouTube auto-captions, of mediocre quality: many formula dictations are garbled, the formulas are reconstructed from the lecture of \cref{ch:c09,ch:c10}
and marked [?] where unsure; the recording ends in the middle of the last sentence).
\end{sources}

\begin{reading}[Where this lecture sits in the other notes]
\emph{Gravity and Light}: lecture 9, ``Newtonian spacetime is curved'' (the same Newtonian gravitational theory, there reformulated: the gravitational field as a component of the connection in a curved Newtonian spacetime, which here is announced and treated in the problem sheet of the previous lecture).
Continues \cref{ch:c10} (interaction fields, force, second axiom, work and potential energy); the next lecture (\cref{ch:c12}) is the two-body problem.
\end{reading}

% ==================================================================
\section{Newton's achievement and the plan}\label{sec:ng-plan}

Good morning, and welcome to the eleventh lecture. Today we look at the gravitational interaction as a special interaction, and as you see, in this lecture we do \emph{not} consider gravitation as curvature of spacetime but quite traditionally as Newton thought of it,
as an interaction (Newton thought of it directly via a force) \ts{11}{1}{0:00}. One must realise what a feat it was of Newton to understand gravitation as a force, and above all as a force which both governs celestial mechanics (the motion of celestial bodies: sun, moon and stars; at least the moon near the Earth)
and \emph{at the same time} the very same law governs how an apple falls from a tree or how one can throw an apple at somebody's head. That these are the same laws was absolutely not clear \ts{11}{1}{0:42}.

\begin{aside}[Aristotle, the aether and the apple (as said in class, not checked)]
If one thinks of antiquity: there was Aristotelian mechanics with all kinds of elements (earth, fire and air, and so on); about the moon Aristotle is said to have said that it must be made of \emph{aether}: it is a natural observation, since the moon is way up there, it must be lighter than air, lighter than earth, water and fire;
if it were heavier it would be down. At the latest with Newton comes the other idea, that there is a universal force of gravitation that governs all this. According to the legend (``it is really the legend'') Newton sat, minutes before the plague broke out in London, under the famous apple tree, and, as it dramatically goes,
the apple fell on his head and this gave him his inspiration; ``I have to think about this'' \ts{11}{1}{1:22}--\ts{11}{1}{2:03}.
(All of this is as said in class and not checked; the details of the legend are not asserted here.)
\end{aside}

We describe it in the modern way, and the context is a \textbf{Newtonian spacetime with absolute space} (``the whole structure, and also an absolute space as we defined it two lectures ago'', \cref{def:as-nrz-as}), on the basis of which the dynamics was formulated last time \ts{11}{1}{2:25}.

\begin{recall}[Particles, interaction fields and the second law (\cref{ch:c10})]
A particle is $(\gamma,m,c_1,\dots,c_L)$: world line, inertial mass, charges; to each charge $c_i$ belongs an interaction field $W_i$; the second Newtonian axiom (\cref{ax:nd-second}) reads
\[
  m\,\nabla_{X_\gamma}X_\gamma\;=\;c_1W_1+\dots+c_LW_L ,
\]
with the $W_i$ evaluated along the world line. (Recall of \ts{11}{1}{2:45}--\ts{11}{1}{3:50}: ``the inertial mass times the acceleration, the covariant derivative of the tangent along its own direction, equals a force built from the interaction fields''.)
\end{recall}
For today we consider only \emph{one} charge, $L=1$, so the indices on the $c$'s may be dropped: $(\gamma,m,c)$ is the particle, $m$ the inertial mass, $c$ the charge \ts{11}{1}{3:06}.

The left-hand side of the second axiom we understand; but it is rather pointless to look for the curve $\gamma$ satisfying the equation if one does not know the right-hand side. Last time I owed you an example of how one \emph{determines} the interaction fields. Therefore we consider today specially the gravitation as an interaction \ts{11}{1}{4:33}.

% ==================================================================
\section{How the interaction fields arise}\label{sec:ng-origin}

The basic idea for all these interaction fields (hinted at last time): \emph{the totality of the particles in the universe (in the Newtonian spacetime with absolute space, ``there stands our basic object'') produces the interaction field}. Each particle contributes to each interaction field, however with a restriction:
\emph{provided that the particle has a corresponding non-vanishing charge} \ts{11}{1}{4:57}.

\begin{intuition}[Passive and active meaning of the charge]\label{int:ng-active}
The charge now has a new meaning, which is not a new notion but an \emph{interpretation}. In the right-hand side of the Newtonian law the charge has a \textbf{passive} meaning: there is already an interaction field produced by all possible particles, and the charge is only there to couple to this field, to \emph{suffer} it.
Here we talk about the \emph{production} of an interaction field by all particles, and it turns out that the very same charges that have the passive meaning have also an \textbf{active} meaning: the totality of the particles produces the interaction field, according to certain \emph{field equations} for such a spatial spacetime vector field \ts{11}{1}{6:02}.
\end{intuition}

\paragraph{The coupled problem.} The general idea: one first solves, if this is at all possible, the field equations, whose result is the field $W$ we need, and puts the result into the second axiom \ts{11}{1}{7:34}. But consider: if one takes a photo, at one absolute time, of the particles in the universe, they produce a certain interaction field; the particles couple to the interaction field produced by their totality and are thereby accelerated and moved;
one second later one takes another photo, with another distribution of particles, and for this new distribution one must solve again the field equations to know what to put into the equation at the next moment. That is the \emph{simplified} view: actually it is a \textbf{coupled problem}, the particles and the fields depend on each other \ts{11}{1}{8:39}.

If the field equations were to be solved at one time and the effect were to be felt \emph{immediately everywhere}, that is not a big problem and one can proceed roughly as just said. But for field equations that do \emph{not} have this property (I will write down some that do) the situation is different: if here an electron moves, or a particle at some time, it takes a time until the field adapts itself everywhere to the new position, and this becomes a very tricky problem \ts{11}{1}{9:47}.

\begin{aside}[Maxwell and radiation (as said in class)]
This occurs historically for the first time in the \emph{Maxwell equations}, and electromagnetic radiation has to do with it: take an electron with its electric field, move it quickly sideways. If the whole world were as one said in school, the electric field would at once look the same, just shifted. But according to the Maxwell equations the electric field does not change so quickly: the particle is moved
and only after some time (with the speed of light) can the field change, and one gets a wave-like effect, the \emph{radiation of energy}; this one learns in the electrodynamics lecture. In Newtonian mechanics we cannot afford this, we could not deal with it; therefore the field equations of Newtonian gravitation are as follows \ts{11}{1}{10:07}--\ts{11}{1}{11:30}.
\end{aside}

\paragraph{The setting: spatial fields and the covector field.} In the Newtonian theory one considers interaction fields (as we saw), and we also said that the interaction fields are all \emph{purely spatial} \ts{11}{1}{11:30}. So, if one takes a spacetime with absolute space, one can draw it again as a foliation: at each time one has from the big metric $G$ a small metric $g$
on each of the spaces at constant time, and on each space one also has the interaction field, its spatial vector field; unlike the metric, however, the interaction field \emph{depends on time}. So: the whole situation can be viewed as an interaction field that in truth is a spacetime field, as a purely spatial field, or as a vector field on the \emph{absolute space} that is, however, time dependent.
This makes the formulae easier, but one could always think about this object in the full spacetime formulation \ts{11}{1}{11:50}.

Then follows a definition that is enormously important. Our spatial vector field has a spatial index $\alpha=1,2,3$ up on its components; from it we define now (``confusingly with the same core letter'') a spatial \emph{covector} field:
\begin{definition}[New concepts: the covector field $W_\alpha$ belonging to $W^\alpha$]\label{def:ng-covector}
Given the interaction field $W^\alpha$ and the spatial metric $g_{\alpha\beta}$ of the absolute space (\cref{def:as-g}), define the \textbf{spatial covector field}
\begin{equation}\label{eq:ng-covector}
  W_\alpha\;:=\;g_{\alpha\beta}\,W^\beta ,
\end{equation}
that is, contract the spatial metric with $W$. It is a covector field not on the whole spacetime manifold but on the absolute space $\mathcal R$ (one can pick an arbitrary one of the spaces and treat it as the absolute space once one is in an inertial system; this is assumed from now on) \ts{11}{1}{13:20}--\ts{11}{1}{14:23}.
\end{definition}
(Later in the lecture a student/the lecturer clarifies the status of this definition, \cref{cn:ng-lowering} below.)

% ==================================================================
\section{The Newtonian gravitational field equations}\label{sec:ng-equations}

Now we formulate the \textbf{Newtonian gravitational field equations}: these are not the second Newtonian axiom but the equations for the interaction field itself. This is the first real section; the rest was preparation and repetition \ts{11}{1}{14:45}. Newtonian gravitation theory: roughly speaking \emph{three} equations, of which only two are real field equations; the third is a condition on the metric
(the metric field is determined by an equation as well) \ts{11}{1}{15:32}.

\begin{definition}[Newtonian gravitational theory: the three field equations]\label{def:ng-equations}
The interaction field $W$ of a Newtonian gravitational interaction (on a Newtonian spacetime with absolute space, in an inertial system, with $\mathcal R$ the absolute space with metric $g$) obeys:
\begin{enumerate}
  \item \textbf{(Curl-freeness)} $(\mathrm dW)_{\alpha\beta}=0$, where for the covector field $W_\alpha$ of \eqref{eq:ng-covector}
        \begin{equation}\label{eq:ng-curl}
          (\mathrm dW)_{\alpha\beta}\;:=\;\frac12\Bigl(\frac{\partial W_\beta}{\partial x^\alpha}-\frac{\partial W_\alpha}{\partial x^\beta}\Bigr)
        \end{equation}
        (the exterior derivative, as on the problem sheet; with the factor $\tfrac12$ of the lecture) \ts{11}{1}{15:32}--\ts{11}{1}{16:36}.
  \item \textbf{(Gauss law, source equation)} Let $D$ be the covariant derivative on the absolute space $\mathcal R$ (the torsion-free metric-compatible one, the Levi-Civita connection of $g$; a big $D$ instead of $\nabla$ so as not to confuse it with the $\nabla$ of the spacetime), whose coefficients are
        \[
          \Gamma^\alpha{}_{\beta\gamma}\;=\;\tfrac12\,g^{\alpha\mu}\bigl(\partial_\beta g_{\mu\gamma}+\partial_\gamma g_{\mu\beta}-\partial_\mu g_{\beta\gamma}\bigr)
        \]
        (\cref{thm:as-christoffel}: ``this was torsion-free and metric-compatible, which we had on the problem sheet''); $D_\alpha$ is shorthand for $D$ in the direction of $\partial/\partial x^\alpha$ in the chart. Then
        \begin{equation}\label{eq:ng-gauss}
          D_\alpha W^\alpha\;=\;-\,4\pi G_{\rm N}\,\rho ,
        \end{equation}
        i.e.\ the contraction of $DW$ (the covariant derivative of the vector field $W$, with the derivative index contracted with the index of $W$) equals $-4\pi G_{\rm N}\rho$; $G_{\rm N}$ is Newton's gravitational constant, $\rho$ a smooth function on the absolute space (continuity would perhaps do), the \textbf{charge density}
        (of the gravitational charge, i.e.\ the density of all particles) \ts{11}{1}{16:36}--\ts{11}{1}{19:46}.
  \item \textbf{(Flatness of the absolute space)} The absolute space is \emph{metrically flat}: there exists an atlas such that in every chart of it
        \begin{equation}\label{eq:ng-flat}
          g_{\alpha\beta}\;=\;\delta_{1\alpha}\delta_{1\beta}+\delta_{2\alpha}\delta_{2\beta}+\delta_{3\alpha}\delta_{3\beta}\;=\;\delta_{\alpha\beta} ,
        \end{equation}
        i.e.\ the matrix of the components of the $(0,2)$ tensor field $g$ has $1,1,1$ on the diagonal and zeros elsewhere \ts{11}{1}{23:58}--\ts{11}{1}{25:22}.
\end{enumerate}
(The sign of \eqref{eq:ng-gauss} was fixed at the end of the lecture, see \cref{cn:ng-sign}; the dictation of the lecture first gives ``$+4\pi G\rho$''.)
\end{definition}

\begin{coursenote}[Lecture errors and corrections on the signs]\label{cn:ng-sign}
In the dictation of the field equation \eqref{eq:ng-gauss} the lecturer first writes $+4\pi G_{\rm N}\rho$; later, while checking the Poisson equation with student interventions, a minus sign is added (\ts{11}{1}{46:50}--\ts{11}{1}{48:34}: ``I want a minus there too''); a few lines later the lecturer first removes the minus below again,
then explains that the correct pair is: $W_\alpha=-\partial_\alpha\Phi$ in the first law and $D_\alpha W^\alpha=-4\pi G_{\rm N}\rho$ in the second, so that the two minus signs combine to the usual $\Delta\Phi=+4\pi G_{\rm N}\rho$ (the Poisson equation) \ts{11}{1}{52:32}. (``One could also leave all minus signs out and it would be consistent, but then the intermediate formulas would
have an unusual sign convention.'') The signs in this chapter are the corrected ones; the captions are partly unreliable here [?]. Another correction: the text of \eqref{eq:ng-curl} has the $\tfrac12$ and the two terms exactly as dictated (``half $\partial W_\beta/\partial x^\alpha$ minus \dots''), but the dictation is garbled and the order of indices is the usual reconstruction [?].
\end{coursenote}

\paragraph{Remarks on the charge density.} The $\rho$ is the charge density of all particles, produced by all particles. Think of the total charge in a region of the absolute space: how does one get it? One chooses a chart $(U,x)$ (if possible covering the region), and then the total charge $C$ (if we consider only one possible charge) is obtained by \emph{integrating} over the chart image: the function $\rho$ lives on the absolute space and must first be transported to the chart, with the inverse chart map,
\begin{equation}\label{eq:ng-totalcharge}
  C\;=\;\int_{x(U)}\mathrm d\alpha^1\,\mathrm d\alpha^2\,\mathrm d\alpha^3\;\sqrt{\det\bigl(g_{\alpha\beta}\bigr)}\;\,\rho\circ x^{-1}\bigl(\alpha^1,\alpha^2,\alpha^3\bigr) .
\end{equation}
Here a factor has to be built in: if one changes charts in the multiple integral, a \emph{Jacobian factor} occurs, and without it the quantity would depend on the chart; the factor is the \emph{square root of the determinant} of the spatial metric, which makes it a true scalar, a chart-independent quantity (exercise on the next sheet: ``childishly easy but one must have computed it once'').
Conversely, to get $\rho$ one takes very, very small regions, notes that the total charge does not change any more, and divides by the volume: ``intuitively you know how this goes'' \ts{11}{1}{19:46}--\ts{11}{1}{23:37}.

\paragraph{The third equation.} In the third equation the \emph{absolute space} is assumed metrically flat. It says: there exists an atlas such that in every chart $g_{\alpha\beta}=\delta_{\alpha\beta}$, so the inner products in these charts simplify: first component times first, plus second times second, plus third times third; the ``standard inner product''. But we recognise that nothing is standard about it; it
depends on a quite special choice of the chart, which must exist at all \ts{11}{1}{25:22}. And Newton says that space is flat; whether this is true we discussed last time (perhaps I fly with a rocket in one direction and after a long time come back from the other side, as Columbus), and it is not really so wrong, as long as we do not extend the domain of physics too much; but we see: we are much better placed,
since we can write this assumption down \emph{explicitly}; it is otherwise always made \emph{implicitly}. And it should change with our ideas about the cosmos (which has already changed a lot) \ts{11}{1}{26:04}. ``Thus these three laws together are the equations that one has to solve: the $g$ occurs implicitly in the first since I lowered the index, and in the second since $D$ is built from the $g$. And so it is damn useful that we have a $g$'' \ts{11}{1}{27:06}.
It becomes incredibly simple, since all the derivatives, if I choose the right atlas (I could choose another atlas, we shall do this later when we use polar coordinates and this beautiful simple form is gone), become simple; but we understand at least what is going on \ts{11}{1}{27:28}.

\begin{theorem}[What Newtonian gravitation theory is]\label{thm:ng-what}
``If someone asks you what Newtonian gravitation theory is, you just recite these three points and that is the answer'': (1) $\mathrm dW=0$, (2) $D_\alpha W^\alpha=-4\pi G_{\rm N}\rho$, (3) the absolute space is metrically flat (and, below, simply connected) \ts{11}{1}{28:08}.
\end{theorem}
(The lecture presents it as the summary; it is a definition of the theory and not a theorem to be proved.)

\subsection{Meaning of the three equations}\label{ssec:ng-meaning}

We extract the meaning of these three requirements, starting with the third \ts{11}{1}{28:29}. In the atlas in which the third holds, all coefficients $\Gamma^\alpha{}_{\beta\gamma}$ of the big $D$ vanish (they are built from derivatives of the constant $g_{\alpha\beta}$): \ts{11}{1}{28:49}
\begin{itemize}
  \item The first equation takes the form $\partial W_\beta/\partial x^\alpha-\partial W_\alpha/\partial x^\beta=0$ with the index lowered; the lowering was defined with the $g$, and the $g$'s, in this chart constants (zeros and ones), may be pulled in front of the derivative; ``the derivative does not feel this $g$ at all, and that is of course a great computational aid'': in another chart, where the $g$'s vary from point to point, one would have to apply the product rule.
        This is just the first law written out carefully \ts{11}{1}{29:10}--\ts{11}{1}{30:52}.
  \item If you were to venture to use a cross-product symbol and do terrible things, then this equation has the same information content as ``the curl of the vector field vanishes''. But the curl as you learned to compute it uses all these special cases (that $g$ is as given), and the notation comes along as if there were no little $g$ in it; but in truth there \emph{is} a little $g$ in this product;
        because one always tells you that it is the standard inner product you do not notice; and this formula generalises neither to other dimensions nor to other inner products. Some people write ``$\operatorname{rot}\vec W=0$'' with the lower-case $g$ up, which works only in three dimensions and is anyway a terrible object. ``But what really goes on is this equation, and perhaps it has a recognition value'' \ts{11}{1}{31:14}--\ts{11}{1}{32:38}.
  \item The second equation takes the form $D_\alpha W^\alpha=\partial_\alpha W^\alpha$ (the $\Gamma$ terms vanish by the third equation) so that (2) becomes the simple form
        \begin{equation}\label{eq:ng-gauss-flat}
          \frac{\partial W^1}{\partial x^1}+\frac{\partial W^2}{\partial x^2}+\frac{\partial W^3}{\partial x^3}\;=\;-4\pi G_{\rm N}\rho ,
        \end{equation}
        the divergence \ts{11}{1}{32:38}--\ts{11}{1}{33:41}.
\end{itemize}
What we have not yet seen is how to use the first equation to rewrite the second; and that is a theorem we need.

\subsection{The potential}\label{ssec:ng-potential}

A first, quick \emph{assumption}, since one has no reason for it at all in general: suppose the interaction field with its index down can be written as the (minus) gradient of a scalar field (smooth on the space):
\begin{equation}\label{eq:ng-potential}
  W_\alpha\;=\;-\,\frac{\partial\Phi}{\partial x^\alpha}\qquad(\Phi\text{ a smooth function on the absolute space}).
\end{equation}
There is in general no reason why a covector field should be representable as minus the gradient of a scalar field. But \emph{if} this is so, then the first equation follows directly: compute
\begin{equation}\label{eq:ng-schwarz}
  (\mathrm dW)_{\alpha\beta}\;=\;\frac12\Bigl(-\frac{\partial}{\partial x^\alpha}\frac{\partial\Phi}{\partial x^\beta}+\frac{\partial}{\partial x^\beta}\frac{\partial\Phi}{\partial x^\alpha}\Bigr)\;=\;0 ,
\end{equation}
since here the same chart is used but different components are inherited, and \emph{Schwarz's theorem} of differential calculus gives zero. The left-hand side is exactly what we called $(\mathrm dW)_{\alpha\beta}$; hence we have found a way to satisfy the first condition: whenever $W_\alpha$ can be written as a gradient of a scalar field the first condition holds \ts{11}{1}{34:46}--\ts{11}{1}{36:52}.
We are interested, however, in the \emph{converse}: what does condition (1) imply? This follows from the Poincar\'e lemma (the name is garbled in the captions: ``Lemma von \dots'', [?]), a result one can look up; it needs a bit more work but is a standard result.

\begin{theorem}[Poincar\'e lemma (converse of the observation) \ts{11}{1}{37:14}]\label{thm:ng-poincare}
If the absolute space (the manifold on which we play the game) is \textbf{simply connected}, then $\mathrm dW=0$ implies that there is a smooth function $\Phi$ on the absolute space with $W_\alpha=-\partial_\alpha\Phi$ (``the same holds for $\mathrm d\Phi$ as we defined it''): the interaction field is a gradient field.
\end{theorem}
(The lecture quotes this without proof; the name of the lemma is not clearly audible, the content is the standard Poincar\'e lemma for closed one-forms on simply connected manifolds.) So for our axioms to work, we must write: \emph{the absolute space is metrically flat and simply connected} \ts{11}{1}{37:58}.

\subsection{Excursion: simply connected}\label{ssec:ng-simply}

What is simply connected? A small excursion, an informal definition \ts{11}{1}{38:23}.
\begin{definition}[Simply connected (einfach zusammenh\"angend), informal; new concept]\label{def:ng-simply}
A topological space (a topological manifold) is \textbf{simply connected} if it is path connected and any continuous closed curve in it can be continuously contracted to a point, as follows. Take a topological space and a point in it, and consider there an arbitrary continuous closed curve $\gamma$ (continuous: preimages of open sets are open). The question: can one extend this to a \emph{family} of continuous curves
$\Gamma\colon\R\times\R\to$ (space), $(\lambda,s)\mapsto\Gamma(\lambda,s)$, again continuous, such that $\gamma$ is a member of the family, $\Gamma(\lambda,0)=\gamma(\lambda)$, and such that the member at the parameter value $1$ is the \emph{trivial curve}, the curve which stays all the time at one point \ts{11}{1}{39:32}--\ts{11}{1}{40:54}?
If such a family exists for every closed curve, the space is simply connected \ts{11}{1}{40:54}.
\end{definition}
Visually: one lays out a lasso in space and can pull it together, without cutting through, until it is contracted to a point; if it works for every lasso, then the space is simply connected \ts{11}{1}{41:15}. Now take a space in which one point is missing, in two dimensions: the plane minus a point. A lasso laid around the missing point can be pulled together but \emph{not to a point}: what is inside is no longer part of the space
accessible to us. It gets stuck, you may wind it five times around, and you do not get it to a fourfold winding either; one can define a winding number; and it depends on the dimension \ts{11}{1}{41:35}. If you are in three-dimensional space and remove a point, the loop can be pulled over the hole (``you know it from bicycle chains: that is an unlucky point, where one locks the bike: if one can pull it off over the top'');
so \emph{three-dimensional space minus a point is simply connected}; but three-dimensional space with an \emph{infinite line} removed is not (one can fasten something to it: the lamp post is practically infinite since one cannot reach the top; ``it is easier than it is'') \ts{11}{1}{42:16}.
It is very important for \emph{the Poincar\'e lemma}, because from the condition $\mathrm dW=0$ on our interaction field one can conclude that this interaction field is already a gradient field \ts{11}{1}{42:58}.

\subsection{The Poisson equation}\label{ssec:ng-poisson}

Back to the subject: with $W$ a gradient field we can use the second equation. Write it with the index of $W$ up; we have learnt $W_\alpha=-\partial_\alpha\Phi$ with the index down; so we must raise the index with the (inverse) metric, $W^\alpha=g^{\alpha\beta}W_\beta$; and since the $g$ may be moved in front of the derivative in the special chart, because it is constant, we get
\begin{equation}\label{eq:ng-poisson-pre}
  -\,g^{\alpha\beta}\,\frac{\partial}{\partial x^\alpha}\frac{\partial}{\partial x^\beta}\,\Phi\;=\;-4\pi G_{\rm N}\,\rho ,
\end{equation}
and in the special chart $g^{\alpha\beta}=\delta^{\alpha\beta}$ the left-hand side is $-\Delta\Phi$, where $\Delta$ is the \textbf{Laplace operator} (you have seen it before; but the Laplace operator depends on the metric $g$ you have, and even this simple form is already not the general one; the general form would contain the $g$'s and the $\Gamma$'s; one can define it without these simplifications) \ts{11}{1}{43:18}--\ts{11}{1}{45:03}. In the chart this is
\begin{equation}\label{eq:ng-poisson}
  \frac{\partial^2\Phi}{(\partial x^1)^2}+\frac{\partial^2\Phi}{(\partial x^2)^2}+\frac{\partial^2\Phi}{(\partial x^3)^2}\;=\;4\pi G_{\rm N}\,\rho ,
\end{equation}
``the second derivative in the 1-direction plus the second derivative in the 2-direction plus in the 3-direction of this field equals $4\pi G_{\rm N}\rho$'': the \textbf{Poisson equation}. This is very, very important, not only for gravitation but also for electrostatics, and one calls it the Poisson equation; but in this form it holds only in a Newtonian spacetime with absolute space where
the absolute space is metrically flat and simply connected and so on and so forth \ts{11}{1}{45:03}--\ts{11}{1}{46:09}. We have so far used all three Newtonian gravitational field equations \ts{11}{1}{46:09}.

\begin{coursenote}[Why the field equations are formulated for $W$ and not for the potential]\label{cn:ng-lowering}
One could begin the other way round, say ``the Poisson equation has to be solved for this scalar field, and once one has it one builds the interaction field by taking the gradient''; but this is not the fundamental statement: the interaction fields are what must be determined, so one must give equations that hold for the interaction field, and not ``ominous equations for ominous scalar functions from which one builds
an interaction field by an ominous rule''. These are the full field equations \ts{11}{1}{48:34}. Also: if you take the Poisson equation $\Delta\Phi=4\pi G_{\rm N}\rho$ and add the corresponding equation for the other interaction (the electric one), you already have the equations of electrostatics, which look exactly like this: ``we have already done two lectures of electrodynamics along the way'' (the same in Newtonian gravitation). ``You will be asked in
various oral exams to write down the Poisson equation or the equation for the gravitational potential'', but we need not these notions; it is this equation \ts{11}{1}{48:55}.
\end{coursenote}

\begin{coursenote}[On ``raising and lowering indices'']\label{cn:ng-lowering2}
(Answer of the lecturer to a student asking about the index lowering; ``that I may not do it'':) The lecturer stresses that nothing is pulled up or down: ``the metric carries your indices'' is bad terminology. One has an object with indices (a vector with upper indices), and by a contraction with the metric one \emph{defines} a \emph{new object} $\omega_\alpha$ with a lower index; by old physics tradition it is also called $W$, though it is another quantity (``musicians write a flat sign here for down'', so one could write $W^\flat$, and with the inverse metric a $W^\sharp$ with the index up again);
since one remembers which is which, the decorations are omitted. So no index is ``pulled''; a new object is defined. In the first law there stands the $W$ with the index down, in the second the $W$ with the index up, ``but we use the same letter, so the terminology is misleading''. If one defines the new object, does one get the old back? Yes: contract with the inverse metric and play the game again, and the old object comes back, \emph{but only because the metric is non-degenerate}: there is no loss of information in lowering and raising,
and we used this when we substituted on the right board, where the index had to be raised again after the first law told us that $W$ with the index down is minus the gradient of a scalar field \ts{11}{1}{49:57}--\ts{11}{1}{52:03}.
\end{coursenote}

% ==================================================================
\section{Solving the field equation}\label{sec:ng-solution}

A second point: we have seen that the whole Newtonian gravitation theory, in the right atlas, \emph{contracts} to the Poisson equation, and we write it down again. But now we look at the \emph{solutions} of the Poisson equation \ts{11}{1}{53:14}. So: in a suitable atlas the three field equations for the gravitational interaction field reduce to $\Delta\Phi=4\pi G_{\rm N}\rho$, with $\rho$ the mass density distribution, the charge density distribution \ts{11}{1}{53:46}.

\begin{theorem}[General solution of the Poisson equation (as said in class: ``there still stands an integral between the general solution and me'')]\label{thm:ng-solution}
Take the field $\Phi$, which lives on the absolute space, but consider its chart representation in a suitable chart $x$ of the suitable atlas (if one is in love with simple notation, one can write the chart representation as a function of the three chart components $\alpha^1,\alpha^2,\alpha^3$, that is, ``the gravitational potential depends on the three space coordinates, namely those of the suitable chart'' as it has always been done).
Then the general solution is
\begin{equation}\label{eq:ng-solution}
  \Phi(\alpha^1,\alpha^2,\alpha^3)\;=\;-\,G_{\rm N}\int\frac{\mathrm d\beta^1\,\mathrm d\beta^2\,\mathrm d\beta^3\;\rho(\beta^1,\beta^2,\beta^3)}{\sqrt{\sum_{i=1}^3\bigl(\alpha^i-\beta^i\bigr)^2}} ,
\end{equation}
where $\rho$ is evaluated at $\beta^1,\beta^2,\beta^3$ and integrated over the relevant region; ``do not be afraid of such triple integrals, you carry them out one after another and, depending on how the region is given, there are tricks; you learn it all, and I make exercises on the problem sheet'' \ts{11}{1}{54:51}--\ts{11}{1}{56:53}. (The lecture first dictates the factor $\tfrac1{4\pi}$ and the $4\pi$ from the field equation cancels, as already settled above, and ``a minus'' in front,
``I hope this is right, this should be correct''; the captions are partly garbled, this is the reconstruction, [?].)
\end{theorem}
(Not proved; the lecture states it.) One might, in exuberance, say for the denominator that this is nothing but the norm of the difference of two position vectors. But you know better by now: the $\alpha^1,\alpha^2,\alpha^3$ are just a triple, and the $\beta^1,\beta^2,\beta^3$ the other triple, which in this suitable chart describe a point; they do \emph{not} form the components of a vector, since there are no position vectors
(we have established this); so also one cannot use the metric to evaluate them; ``nothing of the kind you can do: structurally no inner product is used, structurally no norm is used; it is simply this expression, quite naively, and why does it hold? You can insert it and it comes out.'' It is very tempting to write a norm, but as soon as one begins to compute one has to convert it back; there is no substance in writing a norm there \ts{11}{1}{56:53}--\ts{11}{1}{58:36}.

\begin{example}[The potential outside of a rotationally symmetric body]\label{ex:ng-sphere}
Suppose there is a body whose mass density distribution is rotationally symmetric about a point \emph{with respect to the components of this chart}. Why does this say anything about the real world? Because the chart is so special, with this special choice of $g$, that if one computed the distances of points by computing the lengths of geodesics between them, it would indeed be what in this chart, by this naive measure, is so and so far:
the length of the shortest connections. Hence one can conceive everything like this \ts{11}{1}{58:58}. What is the gravitational potential $\Phi$ outside of a rotationally symmetric body of total charge $C$? You know the expression:
\begin{equation}\label{eq:ng-sphere}
  \Phi(\alpha^1,\alpha^2,\alpha^3)\;=\;-\,\frac{G_{\rm N}\,C}{\sqrt{\sum_{i=1}^3\bigl(\alpha^i-\beta^i\bigr)^2}} ,
\end{equation}
with $C$ the total charge of the body (it comes from integrating $\rho$) and $\beta^1,\beta^2,\beta^3$ the position of the body (its centre) in the chart. Why? One inserts it and computes; ``I do not do it for you, it is simple boring differentiation; it comes on the problem sheet so that you have computed it once'' \ts{11}{1}{1:00:01}--\ts{11}{1}{1:01:04}.
This is the well-known expression for the gravitational potential of a point mass. But this becomes difficult if one wants to introduce point masses: then $\rho$ would have to be written as a delta distribution; these are no longer functions, it is all more complicated. If one takes extended bodies and looks at the potential \emph{outside}, with the centre of mass inside (``I should have said that''), then \eqref{eq:ng-sphere} is the solution; this is also on the problem sheet \ts{11}{1}{1:01:04}--\ts{11}{1}{1:01:46}.
\end{example}
\begin{coursenote}[A student: should there be a $x^{-1}$ and a Jacobian factor?]
A student asks whether it should not be $x^{-1}$ (``you go into the chart'') and whether the Jacobian factor is missing in the integral. The lecturer: ``good question; you integrate here over the space, of course; very good observation: if one is not in this chart, one must put in the Jacobian factor too (the square root of the determinant of $g$ as in \eqref{eq:ng-totalcharge}); but this solution holds only in the chart where we already have the components of this chart; you can compute that the determinant of $g_{\alpha\beta}$ there is $1$, and only if
one changed coordinates would this Jacobian factor play a role. But it is a correct observation: I have taken the simplification already, so that the equations are more familiar. The beginning is important to place it conceptually; for computing, the Poisson equation in this special atlas in this special chart suffices'' \ts{11}{1}{1:02:06}--\ts{11}{1}{1:03:29}.
\end{coursenote}

% ==================================================================
\section{Equivalence of inertial and gravitational charge; the two-particle potential}\label{sec:ng-two}

Then we come to the solution of the \emph{two-body problem}, and there are first several things to remark \ts{11}{1}{1:03:29}. The first remark comes from \textbf{experiment}, explained in one of the first lectures: if in vacuum one takes a small lead ball and a big lead ball and a feather and lets them fall, \emph{they are accelerated equally}, although the lead balls are very massive
and the feather is light \ts{11}{1}{1:03:50}. From this one recognises, by such experiments and a suitable choice of the gravitational constant, that the \emph{gravitational charge must be identical with the inertial mass}. So far the gravitational charge was left open as just some $c$; but it is an experimental observation (``you have perhaps done this in the introductory lab'') that the inertial mass of a particle also fixes its gravitational charge.
This is the \textbf{equivalence of inertial and heavy mass} (Äquivalenz von tr\"ager und schwerer Masse) or, in our terminology, \emph{the inertial mass of a particle also codes its gravitational charge}. This is not something you get from the theory; it is something you \emph{put in} \ts{11}{1}{1:04:11}--\ts{11}{1}{1:05:14}.

\begin{definition}[New concepts: equivalence of inertial and gravitational mass (equivalence principle in this form)]\label{def:ng-equivalence}
For the gravitational interaction one chooses for every particle $(\gamma,m,c)$ the gravitational charge $c=m$ (the inertial mass), after a suitable choice of the gravitational constant \ts{11}{1}{1:04:33}.
\end{definition}
\begin{aside}[Preview]
That the acceleration of a freely falling body is then independent of the mass (in the second axiom both sides contain the same $m$) is the feather experiment; and it is this property which will later make it possible to interpret gravitation as a property of spacetime and not of the particle. (Not said in class here; the lecturer had said in the first lectures that gravitation is not a force; this note only records the link, as announced by him in \cref{ch:c10}.)
\end{aside}

\paragraph{The potential of two particles.} The second remark: when we look at the Poisson equation for \emph{two} particles (we solved it for one particle). Write $\Phi_1$ for the potential of particle $1$ with inertial mass $m_1$ at position $\alpha_{(1)}=(\alpha^1_{(1)},\alpha^2_{(1)},\alpha^3_{(1)})$ in the chart: as we have just seen, it is
\[
  \Phi_1(\alpha^1,\alpha^2,\alpha^3)\;=\;-\,\frac{G_{\rm N}\,m_1}{\sqrt{\sum_{i=1}^3\bigl(\alpha^i-\alpha^i_{(1)}\bigr)^2}} ,
\]
the mass times the gravitational constant divided by this square root; ``only for brevity, really only for brevity, I write the denominator like this: it has no mathematical substance, it is only an abbreviation of the square root of the sum of squares of the differences; it is not a norm'' \ts{11}{1}{1:05:14}--\ts{11}{1}{1:06:36}. But the Poisson equation is \emph{linear} (linear except for the term on the right-hand side, the source), so if we have the solution for one particle we also have the solution for two particles: one simply \emph{adds} the corresponding expression,
\begin{equation}\label{eq:ng-two}
  \Phi(\alpha)\;=\;-\,\frac{G_{\rm N}\,m_1}{\sqrt{\sum_i\bigl(\alpha^i-\alpha^i_{(1)}\bigr)^2}}\;-\;\frac{G_{\rm N}\,m_2}{\sqrt{\sum_i\bigl(\alpha^i-\alpha^i_{(2)}\bigr)^2}} ,
\end{equation}
where $\alpha_{(2)}$ is the position of the second particle; and one computes the gravitational potential at the point $\alpha$, \emph{only at points $\alpha$ that are neither at the position of the first nor at the position of the second particle}, ``so somewhere outside of this first and this second particle'' \ts{11}{1}{1:06:36}--\ts{11}{1}{1:07:47}.

\begin{coursenote}[The recording ends here]\label{cn:ng-end}
The transcript (and the video, length 1:07:47) ends in the middle of the sentence in which the two-particle potential \eqref{eq:ng-two} is discussed; the lecture on the two-body problem as such (the force each particle exerts on the other, Kepler's laws) therefore continues in \cref{ch:c12}. The transcript ends there, so the closing of the lecture (the announcement of the next lecture and of the problem sheet) is not available [?].
\end{coursenote}

% ==================================================================
\section*{Key concepts}
\begin{keyconcepts}
\textbf{Retrieval questions.}
\begin{enumerate}
  \item In what sense is the charge of a particle both passive and active? Why is the problem of motion a coupled problem, and why is it harmless for Newtonian field equations?
  \item Define the covector field $W_\alpha$ belonging to an interaction field. Is the index ``pulled down''?
  \item State the three Newtonian gravitational field equations and the meaning of each ($\mathrm dW=0$, the Gauss equation with $D_\alpha W^\alpha$ and the charge density $\rho$, metrical flatness).
  \item How does one compute the total charge from the density $\rho$ in a chart, and why does $\sqrt{\det g}$ appear?
  \item What is simple connectedness (lasso picture)? What does the Poincar\'e lemma say, and why is it needed to get the potential $\Phi$ with $W_\alpha=-\partial_\alpha\Phi$?
  \item Derive the Poisson equation from the three field equations; where are all three used, and which minus signs combine?
  \item Write the general solution of the Poisson equation and the potential of a rotationally symmetric body; why is $\sqrt{\sum(\alpha^i-\beta^i)^2}$ not the norm of a difference of position vectors?
  \item What is the equivalence of inertial and gravitational mass (experiment, and what is put into the theory)? How does linearity give the potential of two particles?
\end{enumerate}
\textbf{Concepts.}
\begin{itemize}
  \item Active and passive charge (\cref{int:ng-active}); covector field $W_\alpha$ (\cref{def:ng-covector}).
  \item Newtonian gravitational field equations (\cref{def:ng-equations,thm:ng-what}, \eqref{eq:ng-curl}, \eqref{eq:ng-gauss}, \eqref{eq:ng-flat}); charge density and total charge \eqref{eq:ng-totalcharge}.
  \item Potential $\Phi$, Poincar\'e lemma, simple connectedness (\eqref{eq:ng-potential}, \cref{thm:ng-poincare,def:ng-simply}); Poisson equation \eqref{eq:ng-poisson}.
  \item Solution (\cref{thm:ng-solution,ex:ng-sphere}, \eqref{eq:ng-solution}, \eqref{eq:ng-sphere}); equivalence of inertial and gravitational mass and the two-particle potential (\cref{def:ng-equivalence}, \eqref{eq:ng-two}).
\end{itemize}
\end{keyconcepts}

\flushside
\end{paracol}
